% 重构稿：从任务场景到动态图编译，再到协作边上的语义通信架构。
\section{TIDE-Mesh：动态异构拓扑、低熵通信与时态证据记忆}
\label{sec:tide_mesh}

\subsection{应用场景与设计目标}

动态拓扑用于\textbf{任务所需能力、证据来源或处理路径在执行中才逐步明确}的场景。若任务步骤、参与者和验收关系均已确定，固定流水线即可满足需求。对于需要边执行边发现问题的任务，则需根据新证据决定何时增加角色、分出并行调查、加入复核或替换失效成员。

已有系统为这一需求提供了场景参照。Magentic-One通过文件、浏览、编码和执行角色协作处理跨工具任务\cite{tide_magentic_one}；产品架构文档也区分了顺序、并行和动态编排的适用条件\cite{tide_azure_patterns}。本文据此选择三类任务作为设计场景，表~\ref{tab:mesh_scenarios}给出各自的拓扑变化原因。这些场景用于说明机制，不预设本文已经达到相关系统的实测能力。

\begin{table}[htbp]
\centering
\caption{动态协作拓扑的任务场景与触发条件}
\label{tab:mesh_scenarios}
\begin{tabular}{p{0.18\textwidth}p{0.32\textwidth}p{0.39\textwidth}}
\toprule
场景 & 初始工作 & 运行期变化及所需关系 \\
\midrule
软件故障修复 & 定位问题、读取代码和日志 & 根因不明确时并行调查；定位后汇聚给修复角色；变更后调用测试和独立复核；工具失效时切换可行成员 \\
研究报告生成 & 阅读资料、形成证据和论点 & 不同来源并行检索；出现矛盾时增加核验；证据充分后汇聚写作；新增问题触发补充调查 \\
业务资料处理 & 从文档提取需求与业务字段 & 根据资料类型增加专业角色；涉及权限或条款时增加审查；人工变更要求后重建尚未完成的处理分支 \\
\bottomrule
\end{tabular}
\end{table}

上述任务共同要求区分两个决策：\textbf{由谁协作、形成哪些依赖关系；沿这些关系何时通信、传递什么信息}。TIDE-Mesh因此由任务驱动的拓扑控制器和边上的语义通信架构组成。前者输出有任务版本和生命周期的协作图，后者对每条有效边提供一致的消息组织、筛选、压缩及接收校验接口。

\subsection{任务图、活跃协作图与动态流转}

\paragraph{两类图的职责}
任务图$D_t=(\mathcal{T}_t,\mathcal{L}_t)$表达子任务及产物依赖，活跃协作图$G_t=(V_t,E_t)$表达当前由哪些Agent或团队承担子任务，以及哪些角色之间需要交接。候选成员池保持相对稳定，活跃图随当前任务和事件变化。外层规划、质量门和结果汇聚提供控制边界，内层图只在边界内编译和调整。

每个子任务$\tau$记录输入产物、所需能力、硬性资源条件、输出契约及验收条件。控制器先确定任务关系，再分配成员，最后依据生产者与消费者关系生成通信边。每条边可写为
\begin{equation}
e=\langle u,v,\tau,\kappa,v_{\mathrm{task}},b_e,\ell_e\rangle,
\end{equation}
其中$u,v$为发送和接收端，$\kappa$为边类型，$v_{\mathrm{task}}$为任务版本，$b_e$为预算，$\ell_e$为有效期。边类型包括委派、证据供给、产物交接、复核、替补及撤销通知。撤销与重构导致边失效时，应同时处理尚未消费的旧版本消息。

\paragraph{三种基本流线}
初始信息充分且依赖明确时采用顺序交接。存在独立未知因素时，控制器生成并行分支，并在约定的汇聚节点等待必要证据。产物具有较高风险或存在争议时，再引入复核分支。并行范围和复核条件由当前任务决定，不要求每个任务都调用全部成员。

\begin{figure}[htbp]
\centering
\fbox{\parbox{0.92\textwidth}{
\centering
任务状态与新事件 $\rightarrow$ 剩余任务图 $\rightarrow$ 能力互补选择\\[3pt]
$\rightarrow$ 依赖与权限检查 $\rightarrow$ 活跃协作图\\[5pt]
顺序：定位 $\rightarrow$ 修复 $\rightarrow$ 测试\\[3pt]
分支：代码调查 $\parallel$ 环境调查 $\rightarrow$ 证据汇聚\\[3pt]
复核：阶段产物 $\rightarrow$ 独立检查 $\rightarrow$ 通过或修订\\[5pt]
每个箭头均绑定通信契约、版本、预算和生命周期}}
\caption{从任务事件到活跃协作图的编译及典型流线}
\label{fig:mesh_topology}
\end{figure}

\paragraph{软件修复的完整示例}
用户只提供“升级后接口超时”时，初始接入角色先确认环境与复现条件。若无法区分代码变化和部署配置问题，拆出代码调查与环境调查两个分支；二者仅向根因汇聚角色发送相关差异和证据。根因确定后，关闭不再需要的调查边，激活修复到测试的产物交接边。若测试发现涉及数据迁移，再增加迁移检查角色；修复产物满足条件后，由复核角色验收并交付。

任务变更、成员失效和证据冲突分别触发剩余任务更新、失败替补和新增核验边。测试失败时，控制器把反例发回新的修订轮次，而非在同一任务版本内创建无限回环。每个版本内的产物依赖采用可执行的有向无环结构，跨轮反馈通过版本和尝试记录表达；反馈和通知通道可另行存在，但不能绕过调度依赖触发执行。

\paragraph{逻辑协作边与物理通道}
图中的A到B表示B消费A的特定产物，不要求二者建立独立网络连接。单Agent和团队都通过统一通信端口暴露任务输入、消息发送和接收接口；团队内部的交接也可经Supervisor中介。若存在A到B的逻辑边，中介应保留二者标识及边契约，不因中转将其扩大为全员广播。

现有团队内部调用主要经Supervisor组织，Capsule标准链路主要覆盖图节点及其真实后继。本节将相同端口进一步扩展到团队和内部交接边，是目标架构的一部分；需要补齐成员级端口和事件追踪后，才能声称各类协作边均已统一接入。

\subsection{能力互补选择与事件驱动的拓扑优化}

\paragraph{从独立排序转向团队边际贡献}
按单个成员得分取Top-$K$可以快速选择候选，但可能激活多个能力相近的角色，同时遗漏关键能力。本文将选择目标改为：在满足硬约束和并发预算的条件下，覆盖当前任务所需能力，减少重复劳动与跨角色交接开销，并控制频繁切换。

成员画像包括能力、工具、Skill、知识权限、模型配置、历史结果、负载和成本。先用已登记资源和权限构造可行分配关系；某项工具只需由承担对应子任务的成员具备，不能要求团队中每个成员都挂载任务所需的全部工具。未授权、不可用或缺乏必需资源的分配在语义评分前排除。

设待覆盖能力集合为$\mathcal{C}_{\tau}$，成员$i$对能力$c$的适配度为$p_{ic}\in[0,1]$，待选团队为$A$。能力覆盖定义为
\begin{equation}
\operatorname{Cov}(A,\tau)=
\sum_{c\in\mathcal{C}_{\tau}}w_c\max_{i\in A}p_{ic},
\qquad \sum_c w_c=1,
\end{equation}
空集覆盖为零。该定义使已经被充分覆盖的能力具有较小新增收益；必须满足的能力另外设置阈值，不以平均覆盖替代硬性要求。团队目标为
\begin{equation}
\begin{aligned}
F(A,E\mid\tau,S_t)={}&\alpha\operatorname{Cov}(A,\tau)
+\beta\widehat{U}(A,\tau)\\
&-\lambda\sum_{e\in E}\widehat{V}_e
-\mu\operatorname{Red}(A)
-\nu\operatorname{Switch}(A,A_{t-1}),
\end{aligned}
\label{eq:mesh_team_objective}
\end{equation}
其中$\widehat{U}$为预估任务效用，$\widehat{V}_e$为边的预估通信开销，$\operatorname{Red}$为角色职责重复度，$\operatorname{Switch}$为成员切换代价。各项在固定统计口径下归一化，权重由验证集选择。相似角色并不总是冗余：独立复核时，应保留相应角色并检查证据与工具独立性。

\paragraph{可实施的选择算法}
控制器采用有界的边际增益选择：从空团队或可复用成员开始，每次评估加入候选后对覆盖、效用、通信和切换的影响，选择可行且增益最大的成员；硬能力已满足且继续加入无足够收益时停止。随后将成员分配到子任务，依据输入输出契约编译边，再检查依赖、预算、独立复核及并发限制。不满足条件时进行有限修补；仍不可行则返回缺少的能力和资源，不编造成员或强行执行。

G-Designer展示了任务条件化生成协作拓扑的研究方向\cite{tide_gdesigner}，AgentPrune将通信图中的冗余作为优化对象\cite{tide_agentprune}。本文借鉴其优化对象，将显式任务契约保留为执行边界，不直接复用其图生成或剪枝算法。式~\eqref{eq:mesh_team_objective}含有非单调惩罚与依赖约束，有界贪心仅是可解释近似，不声称全局最优或通用近似比。

任务和成员的语义表示可由编码模型产生，候选适配度及预估效用可由配对评分模型学习。启动阶段采用规则、能力目录和带先验的历史统计；积累轨迹后，使用同类任务的验收结果、成本和对照执行训练评分器。仅观察被选成员会形成选择偏差，因此训练集需要受控对照或离线候选评估，不能把未选成员直接标为低质量。图神经网络可作为后续评分器选项，不作为本文架构成立的前提。

\paragraph{事件驱动更新与有界转交}
控制器响应新子任务、能力缺口、成员失败、证据冲突和质量修订，仅更新受影响的剩余子图。能力缺口触发受监督Handoff；运行异常触发Fallback；证据冲突增加核验分支；验收通过后撤销已结束的活跃边。已经验收且输入版本未失效的产物继续复用。

同一尝试中限制连续转交次数和重复参与。若需要让原成员返工，则创建新的修订尝试，携带失败原因与新增证据；不能以禁止任何成员再次参与来阻断合理返工，也不能无限允许同一组合循环。阶段和全局重规划均设置上限。

质量门采用可编码检查与可选模型评审：
\begin{equation}
Q=w_{\mathrm{rule}}Q_{\mathrm{deterministic}}
+(1-w_{\mathrm{rule}})Q_{\mathrm{model}}.
\label{eq:mesh_quality}
\end{equation}
无模型评审时使用独立校准的规则配置，不能把缺失分数当作已获得的模型判断。显式问题未解决时，即使总分较高也不得直接通过。质量门返回通过、修订或重规划，控制器据此更新下一轮任务及活跃关系。

\subsection{低熵通信架构：协作边上的语义传输}

\paragraph{架构定位与跨领域启发}
TIDE-Mesh把通信实现为\textbf{部署在Agent与Agent、Agent与团队、团队与团队逻辑边上的模块链}。对一条$u\rightarrow v$，发送端产物不直接拼入接收者全部历史，而是通过需求编码、胶囊构造、语义选择、保真压缩、接收校验五个模块，再形成$v$的当前输入。模块可以由规则、编码器或模型实现，端口与协议保持一致。

协同感知中的When2com研究通信分组与通信时机\cite{tide_when2com}；Where2comm以空间置信信息选择需要共享的关键区域\cite{tide_where2comm}。本文将这一思路转化为文本协作中的接收者需求和语义单元选择，不使用空间置信图，也不假定视觉网络可直接处理文本。其对应关系是：任务缺口决定何时需要消息，活跃边决定发给谁，验收所需的主张与证据决定发送什么。

“低熵”在本文中表示\textbf{在保持任务必要信息的条件下抑制重复、无关与无依据的信息传播}。它是通信架构的设计名称，不代表最小化字符熵；词项相似度也不等于互信息。系统以通信载荷、字段保真、接收后任务质量和错误采纳率评价效果。

\begin{figure}[htbp]
\centering
\fbox{\parbox{0.92\textwidth}{\centering
发送端Agent/团队 $u$\qquad 接收者契约与已有版本\\[3pt]
$\downarrow$\qquad\qquad\qquad $\downarrow$\\[3pt]
M1 需求编码 $\rightarrow$ M2 结构化胶囊\\[3pt]
$\rightarrow$ M3 受约束语义选择 $\rightarrow$ M4 保真压缩\\[3pt]
$\rightarrow$ M5 版本校验、定向融合 $\rightarrow$ 接收端 $v$\\[3pt]
$\longleftarrow$ 确认、补证请求、冲突与产物质量反馈}}
\caption{部署在逻辑协作边上的五模块语义通信架构}
\label{fig:mesh_communication}
\end{figure}

\paragraph{M1：接收者需求编码}
需求编码器读取边绑定的子任务、接收角色、必需输入、输出契约、可见权限和已确认消息版本，形成接收者需求表示$d_{v,t}$。需求以契约字段为基础，编码模型可补充语义匹配。编码器不访问接收者隐藏推理过程；未知的接收状态按未知处理，不能假定对方已经知道某个前提。

当边上的阶段完成、依赖产物变化、出现约束撤销或接收者请求补证时，通信被触发。普通说明可以合并或等待下一次交接，但执行依赖、取消和冲突通知属于控制信息，不能因预估贡献偏低而静默消失。

\paragraph{M2：结构化胶囊与证据绑定}
胶囊构造器将发送端结果表示为MessageCapsule，将可筛选信息拆为语义单元，并区分主张、证据、约束变化、产物引用及下一动作。每个单元带来源和验证状态，避免压缩后把“待验证假设”改写为“已确认结论”。完整产物留在归档，胶囊可以只传递接收者当前需要的片段和可访问引用。

\begin{table}[htbp]
\centering
\caption{MessageCapsule的通信契约与状态字段}
\label{tab:mesh_capsule_fields}
\begin{tabular}{p{0.25\textwidth}p{0.64\textwidth}}
\toprule
字段类别 & 内容及用途 \\
\midrule
寻址与契约 & 发送端、接收端、边标识、任务及产物版本、消息类型 \\
主张与证据 & 结论、验证状态、证据摘要、来源及原文引用 \\
保护语义 & 约束变化、撤销、失败、冲突和执行阻断条件 \\
增量与基准 & 接收者已确认基础版本、变更字段、删除标记和当前版本 \\
可靠性信息 & 消息标识、校验摘要、有效期、确认状态和压缩血缘 \\
预算与反馈 & 载荷预算、实际用量、补证请求及消费结果 \\
\bottomrule
\end{tabular}
\end{table}

\paragraph{M3：受约束的语义单元选择}
选择器根据接收者当前任务判断信息是否值得传输。对语义单元$z_j$，以任务相关性、证据支持、接收者尚缺信息、时效和重复度估计价值$\widehat{u}_j$。已有词项或向量相似度可用于发现重复，但不能仅据“文本相似”删除数值变更、相反结论或撤销通知。

设$x_j\in\{0,1\}$表示保留单元，$\ell_j$为长度，$\mathcal{P}$为保护单元集合，$\mathcal{K}_v$为接收端必需输入集合，选择问题为
\begin{equation}
\begin{aligned}
\max_{\mathbf{x}}\quad &\sum_j x_j\widehat{u}_j
-\rho\sum_{j<k}x_jx_k\widehat{r}_{jk}\\
\mathrm{s.t.}\quad &\sum_jx_j\ell_j\leq b_e,\qquad
x_j=1\quad(j\in\mathcal{P}),\\
&\operatorname{Cover}(\mathbf{x},\mathcal{K}_v)=1.
\end{aligned}
\label{eq:mesh_semantic_selection}
\end{equation}
$\widehat{r}_{jk}$度量内容重复，覆盖约束要求必要主张连同证据或明确的缺失说明一起交付。若保护语义及必需输入本身超出预算，则该约束集不可行：转为保留完整性的小消息分包、可访问原文引用或显式补证，不悄悄丢弃保护字段。无法满足接收前提时，接收者等待或暂停。

工程路径可采用保护单元优先、按边际价值和长度选择其余单元的近似算法。预估价值可以由轻量模型学习，但必须使用固定任务质量口径校准。新颖性只描述接收者是否见过相关信息，不能证明其正确。

\paragraph{M4：字段感知压缩与保真检查}
压缩器对被选中的说明文字去重、摘要和缩短，优先保留事实值、单位、时间、约束变更、验证状态及来源关系。模型可以压缩描述字段，符号检查器核对受保护字段及必需引用；压缩失败时回退到未改写的结构化字段或缩小普通说明范围。

字段一致不保证所有自然语言含义均一致。因此对高风险结论还需执行来源核验或独立复核，并记录压缩前后内容及父消息标识。完整原文在归档中保留，使接收者能按需展开证据，而不是依赖不可追溯的短摘要。

\paragraph{M5：版本校验、定向融合与反馈}
接收模块检查权限、任务版本、有效期、消息标识和校验摘要；重复消息不重复提交业务状态。增量必须相对\textbf{该接收者已确认的基准}计算。若发送端上一次消息被筛除或未确认，不能继续以其发送版本作为接收端基础版本。

设发送者$u$对接收者$v$的已确认视图为$\widehat{S}_{u,v}^{b}$，待发送视图为$S_u^{k}$，则
\begin{equation}
\Delta_{u,v}^{b\rightarrow k}
=\operatorname{Diff}(\widehat{S}_{u,v}^{b},S_u^{k}).
\end{equation}
接收者校验基础版本后合并增量并返回确认；缺少基础版本时请求完整快照或缺失增量，不能直接套用。删除标记只作用于该业务视图；全局共享状态仍按第\ref{sec:tide_anchor}节的版本提交规则更新。

确认投递和确认已消费分别记录。运行恢复时同时核对信箱、确认进度和基础视图；在可能重复投递的条件下，通过消息去重和提交幂等性避免重复生效，不宣称普遍的端到端恰好一次执行。控制通知过期后应重新获取当前有效状态，不能因TTL到期就继续使用旧约束。

校验后的语义内容进入接收者当前输入，协议元数据保留在运行记录中。接收者可以返回已足够、缺少证据、版本冲突或任务不再需要等反馈；来源正确性仍需独立检查，接收确认不等于事实已验证。

\subsection{时态证据记忆与跨模块一致性}

时态证据层记录“哪个来源在什么时间、基于什么证据提出了哪个版本的事实”。通信胶囊先进入消息记录或候选事实区，只有完成相应验收后，才按第\ref{sec:tide_anchor}节的准入规则形成TASK、PROJECT或GLOBAL记忆。这样，消息到达、共享状态提交和经验入库成为三个有区别的事件。

事实记录包含实体、属性、取值、有效时间、记录时间、证据和验证状态。同一槽位在不兼容的时间范围内取不同值时，保留冲突供核验；明确的新版本可以替代旧版本，但历史不得被无痕覆盖。自然语言矛盾不能只靠字符串比较判定。

当事实或记忆被撤销时，系统沿已记录的消费依赖向受影响接收者发送保护通知，使其停止使用相关版本，并由共享状态管理器决定补证或局部重执行。通信模块负责可靠传递撤销信息，TIDE-Anchor负责管理事实和产物依赖，拓扑控制器负责为修复任务重新组织成员，避免多个模块分别决定同一产物的权威状态。

\subsection{核心算法、实现边界与评价方法}

\begin{algorithm}[htbp]
\caption{任务契约驱动的活跃协作图编译}
\label{alg:mesh_collaboration}
\KwIn{共享状态$S_t$，事件$e_t$，候选成员池$\mathcal{A}$，预算$b$}
\KwOut{可执行协作图或不可行原因}
更新受影响的剩余任务，确定能力、输入输出和验收条件\;
按资源、权限和可用性过滤不合法的任务--成员分配\;
以边际贡献选择互补成员，检查必需能力与并发限制\;
根据产物依赖生成交接、汇聚和复核边，绑定通信契约\;
检查版本内无环、资源预算与复核独立性，执行有限修补\;
\If{仍不可行}{\Return 缺少的能力、资源或契约条件\;}
执行就绪子任务，接收质量、失败和能力缺口事件\;
关闭结束或失效边；在尝试预算内更新剩余子图\;
\Return 图版本及运行记录\;
\end{algorithm}

\begin{algorithm}[htbp]
\caption{协作边上的接收者条件化语义通信}
\label{alg:mesh_transport}
\KwIn{发送端结果$y_u$，有效边$e=(u,v)$，已确认视图$B_{u,v}$}
\KwOut{接收结果、补证请求或阻断原因}
M1：读取接收者契约、缺失输入和已确认版本\;
M2：构造带来源、验证状态和保护语义的胶囊\;
M3：在预算内选择语义单元，保留必需信息与依赖\;
\If{必要字段无法装入预算}{采用分包或引用；无法完整交付则阻断\;}
M4：压缩普通描述，检查保护字段和引用，失败则回退\;
M5：检查权限、有效版本及消息去重；核对增量基础版本\;
\eIf{基础版本不一致}{请求快照或缺失增量\;}{
    定向写入信箱，按契约融合，记录确认及消费状态\;
}
反馈信息缺口、冲突和验收结果，更新接收者已确认视图\;
\end{algorithm}

设候选数为$N$，最大激活成员数为$K$，能力维度为$d$，一次候选效用及局部边估计成本为$T_s$。直接实现的边际选择成本约为$O(KN(d+K+T_s))$；任务图编译后的依赖检查为$O(|V|+|E|)$，有限图修补需单独计入。该方法比单独排序开销更高，其价值应由任务质量和总执行成本检验。

设单条消息具有$J$个语义单元、预算离散值为$B$。不含成对重复项的背包动态规划成本为$O(JB)$；完整式~\eqref{eq:mesh_semantic_selection}含成对项和覆盖约束，不能直接套用这一复杂度。工程近似若缓存成对相似度，其预处理为$O(J^2T_s')$，其中$T_s'$为单次相似度成本；接收者数和核验模型调用另计。总载荷按实际投递累计：
\begin{equation}
V_{\mathrm{comm}}=\sum_{m\in\mathrm{Delivered}}
\operatorname{Size}(m).
\end{equation}
分包、确认、补证、重传和快照都应计入相应通信统计，不以潜在边数直接替代实际Token开销。

现有工程提供资源硬过滤、多指标Top-$K$、受监督转交与替补，以及图节点边界的胶囊门控、压缩和状态增量。能力互补的图编译、统一成员端口、语义单元选择、接收者确认基准和撤销传播构成本节扩展设计，须在后端接入后验证。原有规则路径可以作为对照和冷启动回退，不据其测试结果宣称扩展设计已经通过实验。

拓扑实验应对比固定流水线、固定团队、原Top-$K$与能力互补图编译，分别观察任务完成、能力缺口、切换与编排耗时。通信实验在相同任务和拓扑下，对比原文传递、结构化胶囊、语义选择和完整五模块链，报告含模块调用开销的Token、时延、证据保留、压缩失真及错误采纳率。丢包、版本错位、重复投递和撤销通知用于验证协议行为，不能以消息变短直接推断任务质量提高。

